\chapter{The Risk-Management Function}\label{ch:fm:the-risk-management-function}

The examiner appointed in the bankruptcy of Lehman Brothers reported in March 2010 that between December 2006 and December 2007 the firm raised its
firm-wide \omterm{def:fm:the-head-of-desks-job:appetite}{risk appetite} limit three times, from \$2.3 billion to \$4.0 billion. The last increase was approved in January 2008 and backdated to
December 2007, which removed the excesses of the intervening weeks from the record; the examiner's advisers calculated that the method used the year
before would have given a limit of about \$2.5 billion. The firm's stress tests left out its commercial real estate and private equity investments
and, for a time, its leveraged loans. The numbers were there. What failed was a function that could not make them count.

\section{Independence: where the risk function reports}

A trading firm's risk function measures the risks its businesses take, sets and polices the limits of chapter 7, and objects when the two disagree.
It is the second of Book 6's three lines of defence: the businesses own their risks, the risk function challenges them, and internal audit checks
both. Its value lies entirely in the objection, and an objection is only as strong as the reporting line behind it.

\begin{definition}[Chief risk officer, risk committee]\label{def:fm:the-risk-management-function:cro}
The \emph{chief risk officer}\index{chief risk officer} is the executive responsible for the firm's risk function: for setting the risk limits across
the firm and monitoring compliance with them, for the policies and controls that implement the \omterm{def:fm:the-head-of-desks-job:appetite}{risk appetite}, and for reporting breaches, deficiencies
and emerging risks. The \emph{risk committee}\index{risk committee} is the body, of the board or of senior management, that approves the risk
policies and limits, receives the chief risk officer's reports, and decides on the exceptions and increases the policies reserve to it.
\end{definition}

Three things make the \omterm{def:fm:the-risk-management-function:cro}{chief risk officer} independent in practice: a reporting line that does not run through the businesses whose risk is measured,
pay that does not depend on their revenue, and a direct line to the \omterm{def:fm:the-risk-management-function:cro}{risk committee} that no one in between can filter. The first is visible on an
organisation chart (\cref{fig:fm:the-risk-management-function:org}); the others are not, and matter as much. The examiner's report cites press
reports, published before the bankruptcy, that the firm had removed its \omterm{def:fm:the-risk-management-function:cro}{chief risk officer} in 2007 because of opposition to its growing
accumulation of illiquid investments; the \omterm{def:fm:the-risk-management-function:cro}{chief risk officer} had also argued for a lower limit increase at the end of 2006 than the one adopted.

\begin{omfigure}
\begin{tikzpicture}[font=\footnotesize,>=stealth,box/.style={draw,rounded corners,align=center,fill=omDef!10,minimum height=0.75cm,text width=2.7cm}]
\node[box,fill=gray!15] (board) at (0,3.3) {board};
\node[box,fill=omProp!15] (rc) at (-3.4,2.1) {risk committee};
\node[box,fill=gray!15] (ceo) at (3.4,2.1) {chief executive};
\node[box,fill=omProp!15] (cro) at (-3.4,0.6) {chief risk officer};
\node[box] (bus) at (3.4,0.6) {head of markets};
\node[box] (mr) at (-5.0,-0.9) {market and credit risk};
\node[box] (mv) at (-1.9,-0.9) {model validation};
\node[box] (d1) at (1.9,-0.9) {desks};
\node[box] (d2) at (5.0,-0.9) {trading support};
\draw[->] (board) -- (rc); \draw[->] (board) -- (ceo);
\draw[->,very thick,omProp] (rc) -- (cro); \draw[->,very thick,omProp] (ceo) -- (cro);
\draw[->] (ceo) -- (bus); \draw[->] (cro) -- (mr); \draw[->] (cro) -- (mv); \draw[->] (bus) -- (d1); \draw[->] (bus) -- (d2);
\draw[<->,dashed,gray] (cro) -- node[above,font=\scriptsize,text=black] {challenge} (bus);
\end{tikzpicture}

\omcaption{An independent risk function on an organisation chart: the \omterm{def:fm:the-risk-management-function:cro}{chief risk officer} reports to both the \omterm{def:fm:the-risk-management-function:cro}{risk committee} and the chief executive
(thick lines), not to the head of the businesses whose risk is measured, and challenges them from the side. Schematic.}\label{fig:fm:the-risk-management-function:org}
\end{omfigure}

\begin{dated}{2026-09}{The chief risk officer in US bank rules}\label{dat:fm:the-risk-management-function:rules}
Under the Federal Reserve's Regulation YY (12 CFR 252.33), a US bank holding company with average total consolidated assets of \$100 billion or more
must maintain a \omterm{def:fm:the-risk-management-function:cro}{risk committee} of its board: an independent committee with a written charter, chaired by an independent director, meeting at least
quarterly and documenting its decisions. It must appoint a \omterm{def:fm:the-risk-management-function:cro}{chief risk officer} experienced in the risks of large, complex financial firms, responsible
for establishing enterprise-wide risk limits and monitoring compliance with them, whose pay must be consistent with an objective assessment of the
risks taken, and who reports directly to both the \omterm{def:fm:the-risk-management-function:cro}{risk committee} and the chief executive officer. Other jurisdictions and smaller firms set their
own rules; a proprietary trading firm may have no regulatory requirement at all.
\end{dated}

Supervisors who compared major firms after the first months of the 2007 turmoil found the difference was not independence alone. The Senior
Supervisors Group reported in March 2008 that at firms which avoided significant losses the risk function had independence and authority but also
considerable direct interaction with senior business managers; that some firms escalated concerns about emerging risks as early as the summer of
2006, gaining up to a year; and that in others hierarchical structures filtered what reached senior management. A risk function that is independent
and remote sees the numbers late and is heard little.

\section{Committees and the risk officer's tools}

The risk function's tools are few and blunt. It proposes the \omterm{def:fm:the-head-of-desks-job:appetite}{risk appetite} (chapter 6) and the \omterm{def:fm:limits-and-capital-allocation:framework}{limit framework} (chapter 7); it measures value at
risk and runs stress tests (Book 6); its model validation team approves the models that produce those numbers; it signs off new products and new
strategies; and it decides, or recommends, what happens when a limit is breached. Each tool works only if a trade cannot happen without it.

\begin{definition}[Four-eyes principle]\label{def:fm:the-risk-management-function:foureyes}
The \emph{four-eyes principle}\index{four-eyes principle} requires that a decision or action of a defined kind (a limit increase, a model change, a
new product, a payment) be approved by a second person, independent of the one who proposes or executes it, before it takes effect.
\end{definition}

The principle is the operational form of segregation of duties (Book 6): the trader who wants a larger limit does not grant it, and the quant who
changes a risk model does not validate the change. It fails when the second pair of eyes belongs to someone whose interests are the first person's:
a head of desk who approves increases for their own desk applies four eyes on paper and two in fact.

\begin{method}[Running a risk committee]\label{met:fm:the-risk-management-function:committee}
\begin{enumerate}
\item Write its charter: members, quorum, what it alone may approve (increases above a threshold, new products, model changes that move a limit
measure), and how often it meets.
\item Give it standing reports from the risk function, not from the businesses: utilisation against limits, breaches and their resolution, stress
results, exceptions outstanding.
\item Record every decision with its reason and its expiry; review expired exceptions at every meeting.
\item Give the \omterm{def:fm:the-risk-management-function:cro}{chief risk officer} the right to escalate to the board's committee directly, and use it.
\end{enumerate}
\end{method}

\section{Escalation: breaches, excesses and exceptions}

Chapter 7 defined the \omterm{def:fm:limits-and-capital-allocation:breach}{limit breach} and the difference between hard and \omterm{def:fm:limits-and-capital-allocation:breach}{soft limits}. What happens next is the \omterm{def:fm:the-risk-management-function:escalation}{escalation procedure}.

\begin{definition}[Escalation procedure, risk exception]\label{def:fm:the-risk-management-function:escalation}
An \emph{escalation procedure}\index{escalation procedure} states, for each severity of \omterm{def:fm:limits-and-capital-allocation:breach}{limit breach} or control failure, who must be told, within
what time, who may approve a temporary increase or other exception, and what must be done if no approval is given. A \emph{risk exception}\index{risk
exception} is a documented, approved and time-limited departure from a limit or policy, with an owner and an expiry date.
\end{definition}

A breach is graded by size and by duration. The chapter's rules (\cref{lst:fm:the-risk-management-function:rules}) grade a breach severity 1 if
it exceeds the \omterm{def:fm:limits-and-capital-allocation:breach}{hard limit} by at most 5\% for at most two days, severity 3 if it exceeds it by more than 20\% or lasts more than five days, and
severity 2 in between. Severity 1 is reported to the head of desk the same day, severity 2 to the \omterm{def:fm:the-risk-management-function:cro}{chief risk officer} within a day, and severity 3 to
the \omterm{def:fm:the-risk-management-function:cro}{risk committee} within two. There are three honest ways out of a breach: cut the position, have a temporary increase approved by someone
independent, or show that the risk was mismeasured and have a corrected model validated. The same three are also the dishonest ways out, when the
increase is granted to fit the position or the model is changed to fit the limit.

\omcode{code/firm/escalation/firm_escalation.py}{28}{53}{The escalation rules as data: severity by size and duration, and who is told within
how many business days.}\label{lst:fm:the-risk-management-function:rules}

\begin{omfigure}
\begin{tikzpicture}
\begin{axis}[width=10.5cm,height=5.2cm,xlabel={\small month},ylabel={\small \$ billion},xmin=0,xmax=15,ymin=0,ymax=5,
  xtick={0,4,7,10,13,15},xticklabels={Nov 06,Mar 07,Jun 07,Sep 07,Dec 07,Feb 08},
  tick label style={font=\footnotesize},grid=major,grid style={gray!20},table/col sep=comma,
  legend cell align=left,legend style={font=\scriptsize,at={(0.5,-0.3)},anchor=north,/tikz/every even column/.append style={column sep=8pt}},legend columns=2]
\addplot[omDef,very thick,const plot] table[x=month,y=limit]{figdata/desk/12-the-risk-management-function/lehman.csv};
\addplot[omProp,thick,dashed] coordinates {(13,2.5) (15,2.5)};
\node[font=\scriptsize,align=right,anchor=south east] at (axis cs:12.8,4.05) {approved Jan.\ 2008,\\backdated to 3 Dec.\ 2007};
\node[font=\scriptsize,align=left,anchor=north west] at (axis cs:1.2,3.2) {chief risk officer argued\\for \$2.6--2.7 billion};
\legend{firm-wide risk appetite limit,2008 limit on the 2007 method}
\end{axis}
\end{tikzpicture}

\omcaption{The firm-wide \omterm{def:fm:the-head-of-desks-job:appetite}{risk appetite} limit in the Lehman examiner's report: \$2.3 billion, raised to \$3.3 billion for fiscal 2007, to \$3.5
billion on 7 September 2007 (the firm was over the new limit on every business day of September but one) and to \$4.0 billion by an approval of 14
January 2008 backdated to 3 December 2007; the examiner's advisers put the 2008 limit on the previous year's method at about \$2.5 billion. Source:
the chapter's ledger.}\label{fig:fm:the-risk-management-function:lehman}
\end{omfigure}

The public record (\cref{fig:fm:the-risk-management-function:lehman}) shows each of the dishonest exits. A limit raised while the firm was over
it; an increase backdated so that the breach disappeared from the register; a limit recalculated with new assumptions that produced a larger number;
and, before these, a large position left out of the usage calculation from May to August 2007, which kept the usage under the limit until it was
included. One Quant Book 6 tells the story of a later case in which a risk model was changed while a limit was being breached. The escalation
register must be built so that such patterns are visible: limits and models kept point in time, as they were known each day, and every change
recorded with the day it was approved as well as the day it takes effect.

\begin{proposition}[Waiting for the committee]\label{prop:fm:the-risk-management-function:wait}
If the \omterm{def:fm:the-risk-management-function:cro}{risk committee} meets every $m$-th business day and a breach opens on a day uniformly distributed over the cycle, the wait until the next
meeting after the opening day is uniform on $\{1,\dots,m\}$, with mean $(m+1)/2$ business days.
\end{proposition}

\begin{proof}
The next meeting after day $t$ is day $t+1+((-(t+1))\bmod m)$; as $t$ runs over a cycle, the wait $1+((-(t+1))\bmod m)$ takes each value
$1,\dots,m$ once.
\end{proof}

With weekly meetings the wait is three business days on average. A procedure that makes the desk wait with its breach open lengthens every breach
that goes to the committee; one that makes the desk cut while it waits does not.

\section{Tutorial: a year of breaches}\label{tut:fm:the-risk-management-function:tut}

\begin{tutorial}
\textbf{Goal.} Run the escalation rules over a year of \omterm{def:fm:limits-and-capital-allocation:breach}{limit utilisation} for twenty desks, measure how breaches are resolved and how long they last,
and flag the pattern of the public record. \textbf{End state:} a table and two charts of how breaches end (\cref{tab:fm:the-risk-management-function:regimes},
\cref{fig:fm:the-risk-management-function:regimes} and \cref{fig:fm:the-risk-management-function:window}).
\begin{tutsteps}
\item \textbf{The year.} \texttt{fm\_risk.simulate(regime)} draws each desk's value at risk against a \omterm{def:fm:limits-and-capital-allocation:breach}{hard limit} of 100 (a persistent random walk
around 75\% utilisation). When a hard breach opens the desk asks for a temporary increase (probability 0.45), changes its risk model (0.08; the
measured risk falls by a fifth) or cuts its position.
\item \textbf{Three regimes.} The head of desk approves increases the next day; or the \omterm{def:fm:the-risk-management-function:cro}{risk committee}, meeting every fifth day, approves them with
probability 0.6 and validates model changes at the meeting; or the same committee, with the desk required to cut while it waits.
\item \textbf{The register.} \texttt{firm.escalation.register} finds every breach, grades it, classifies how it closed and flags limits raised or
models changed during a breach and repeated increases (\cref{lst:fm:the-risk-management-function:register}).
\item \textbf{The comparison.} \texttt{fm\_risk.compare()} pools ten simulated years for each regime.
\end{tutsteps}
\end{tutorial}

\omcode{code/firm/escalation/firm_escalation.py}{77}{111}{The breach register: every run of days above the hard limit, graded, classified by
how it closed and checked for red flags.}\label{lst:fm:the-risk-management-function:register}

\begin{table}[ht]
\centering\small
\begin{tabular}{lrrr}
\toprule
 & head of desk & \omterm{def:fm:the-risk-management-function:cro}{risk committee} & committee, \\
 & approves & approves & cut while waiting \\
\midrule
breaches a year (20 desks) & 27.8 & 26.4 & 27.8 \\
closed by a position cut (\%) & 50.0 & 78.4 & 87.4 \\
closed by a limit increase (\%) & 45.3 & 14.4 & 5.4 \\
closed by a model change (\%) & 3.2 & 6.1 & 6.5 \\
mean duration (days) & 1.15 & 1.64 & 1.25 \\
breach-days a year & 31.8 & 43.5 & 34.6 \\
severity 2 or 3 (ten years) & 32 & 57 & 25 \\
desk-days above the original limit (\%) & 2.44 & 1.33 & 0.74 \\
breaches flagged (\%) & 54.7 & 23.1 & 11.9 \\
\bottomrule
\end{tabular}
\caption{Ten simulated years of breaches on twenty desks under three approval regimes (0.7--1.4\% of breaches are still open at a year end).
Data: \texttt{fm\_risk.compare}.}\label{tab:fm:the-risk-management-function:regimes}
\end{table}

Most breaches are small and short: the median lasts one day in every regime, since utilisation that drifts over a limit often drifts back.
The regimes differ in how the rest end (\cref{tab:fm:the-risk-management-function:regimes}). When the head of desk approves increases, 45\% of
breaches are closed by raising the limit; breaches are the shortest, and the desks spend 2.44\% of their days above the limit they started the year
with. Requiring the committee's approval cuts increases to 14\% and time above the original limit to 1.33\%, but breaches last longer while
desks wait for a meeting (\cref{prop:fm:the-risk-management-function:wait}): 43.5 breach-days a year against 31.8, and more of them reach
severity 2 or 3. Making desks cut while they wait keeps the committee's control and removes the delay: 5\% of breaches end with an increase,
the mean duration is 1.25 days, and the desks spend 0.74\% of their days above their original limits.

\begin{omfigure}
\begin{tikzpicture}
\begin{axis}[width=10.5cm,height=4.6cm,xbar stacked,bar width=12pt,xmin=0,xmax=100,xlabel={\small share of breaches (\%)},
  ytick={0,1,2},yticklabels={head of desk approves,risk committee approves,{committee, cut while waiting}},y dir=reverse,ymin=-0.6,ymax=2.6,
  tick label style={font=\footnotesize},grid=major,grid style={gray!20},table/col sep=comma,
  legend cell align=left,legend style={font=\scriptsize,at={(0.5,-0.35)},anchor=north,/tikz/every even column/.append style={column sep=8pt}},legend columns=4]
\addplot[fill=omDef!55,draw=omDef] table[y=pos,x=cut]{figdata/desk/12-the-risk-management-function/regimes.csv};
\addplot[fill=omProp!55,draw=omProp] table[y=pos,x=increase]{figdata/desk/12-the-risk-management-function/regimes.csv};
\addplot[fill=omThm!50,draw=omThm] table[y=pos,x=model]{figdata/desk/12-the-risk-management-function/regimes.csv};
\addplot[fill=gray!40,draw=gray] table[y=pos,x=open]{figdata/desk/12-the-risk-management-function/regimes.csv};
\legend{position cut,limit increase,model change,still open}
\end{axis}
\end{tikzpicture}

\omcaption{How breaches end under three approval regimes: ten simulated years of twenty desks each. Moving the approval of increases away from
the desk moves breaches from the limit to the position. Data: \texttt{fm\_risk.compare}.}\label{fig:fm:the-risk-management-function:regimes}
\end{omfigure}

\begin{omfigure}
\begin{tikzpicture}
\begin{axis}[width=10.5cm,height=5cm,xlabel={\small business days from the model change},ylabel={\small share of the original limit},xmin=-15,xmax=24,
  ymin=0.6,ymax=1.25,tick label style={font=\footnotesize},grid=major,grid style={gray!20},table/col sep=comma,
  legend cell align=left,legend style={font=\scriptsize,at={(0.5,-0.3)},anchor=north,/tikz/every even column/.append style={column sep=8pt}},legend columns=2]
\addplot[omDef,thick,mark=*,mark size=1pt] table[x=day,y=utilisation]{figdata/desk/12-the-risk-management-function/window.csv};
\addplot[omProp,thick,const plot] table[x=day,y=limit]{figdata/desk/12-the-risk-management-function/window.csv};
\draw[gray,dotted] (axis cs:0,0.6) -- (axis cs:0,1.25);
\legend{measured value at risk,hard limit (with temporary increases)}
\end{axis}
\end{tikzpicture}

\omcaption{One desk around the first breach of the simulated year closed by a model change: a temporary increase to 110, a one-day lapse when it expires and a new
increase, then a breach of 7.7\% closed two days later by a model change that cuts the measured risk by a fifth with no change in the position. The
register flags both the repeated increases and the model change. Data: \texttt{fm\_risk.model\_change\_window}.}\label{fig:fm:the-risk-management-function:window}
\end{omfigure}

\textbf{What to change next.} Record each increase with its approval day and check it with \texttt{firm.escalation.backdated}; give the desks a
position that is left out of the usage for three months and see which register metric notices it first.

\section{What the risk officer looks for}

A register that only counts breaches misses the failures of the public record, which were patterns rather than events. The risk officer looks for:
\begin{itemize}
\item limits raised, or models changed, while a breach is open, and temporary increases that are renewed rather than allowed to expire;
\item changes that take effect before they were approved;
\item positions outside the measure: exposures left out of the usage calculation or out of the stress tests because they are new, illiquid or
hard to model (the examiner found the firm's stress tests excluded its commercial real estate and private equity investments);
\item limits recalculated with new assumptions that produce a larger number at the moment a larger number is needed;
\item exceptions without an owner or an expiry, and a committee that approves everything it is asked.
\end{itemize}

\begin{definition}[Risk culture]\label{def:fm:the-risk-management-function:culture}
\emph{Risk culture}\index{risk culture} is the set of norms, attitudes and behaviours in a firm that decide how risk is taken and controlled in
practice: whether limits are treated as binding, whether bad news travels up quickly, and whether those who object are heard.
\end{definition}

Culture cannot be read off a register, but its traces can: the share of breaches closed by increases, the share of flagged breaches, the time
from a breach to the committee's hearing of it, the number of exceptions past their expiry. The examiner found that management had treated the
firm-wide limit as a ``soft'' guideline while describing it to its regulator and board as a meaningful constraint. On the question the examiner
had to answer, the report found insufficient evidence that the senior officers' conduct in managing risk fell outside the business judgement rule or
was reckless or irrational: decisions that proved unwise were, legally, still decisions management was entitled to make. The lesson for a risk
function is that its authority must be written into the procedure, since it cannot rely on the law to supply it after the event.

\begin{remark}[Chapter 7's limits and this chapter's procedure]\label{rem:fm:the-risk-management-function:ch7}
Chapter 7 sized limits and priced capital as if every limit held. This chapter's register is how a firm finds out whether they do: a limit that is
raised whenever it binds is not a limit, and the capital it was meant to protect is not protected.
\end{remark}

\section{Build: the escalation register}\label{bld:fm:the-risk-management-function:escalation}

\begin{build}
\bfield{Purpose} The breach register and escalation workflow as data: grading, notification, resolution tracking, audit metrics and red flags.
\bfield{Interface} \texttt{firm.escalation}: \texttt{Rules}, \texttt{severity}, \texttt{notify}, \texttt{hard\_path} (from a
\texttt{firm.limitalloc} node and dated increases), \texttt{register}, \texttt{stats}, \texttt{flagged}, \texttt{backdated}.
\bfield{Rules} A breach is a run of days above the effective \omterm{def:fm:limits-and-capital-allocation:breach}{hard limit}; it closes by a limit increase, a model change or a position cut, in that
order of precedence; flags are recorded, never cleared; limits and models are point in time.
\bfield{Acceptance tests} \texttt{code/firm/escalation/tests/}: the severity grid; one breach of each resolution type on hand-made data; repeated
increases; a backdated change.
\bfield{Stretch} Notification deadlines checked against a log of who was told when; exceptions with owners and expiry dates; a daily report for the
\omterm{def:fm:the-risk-management-function:cro}{risk committee}.
\end{build}

\begin{omsources}
\item Report of Anton R. Valukas, Examiner, In re Lehman Brothers Holdings Inc., No.~08-13555 (Bankr. S.D.N.Y.), 11 March 2010, volume 1.
\item 12 CFR 252.31 and 252.33 (Regulation YY).
\item Senior Supervisors Group, \emph{Observations on Risk Management Practices during the Recent Market Turbulence}, 6 March 2008.
\end{omsources}

\section{Exercises}

\begin{exercise}[$\star$]\label{exo:fm:the-risk-management-function:1}
Grade under the chapter's rules a breach of 3\% lasting one day, one of 10\% lasting one day and one of 3\% lasting six days, and say who is told and
by when.
\end{exercise}

\begin{exercise}[$\star$]\label{exo:fm:the-risk-management-function:2}
By what percentage did the firm-wide limit of \cref{fig:fm:the-risk-management-function:lehman} rise from November 2006 to December 2007? By how much
did the 2008 limit exceed the one the previous year's method gave?
\end{exercise}

\begin{exercise}[$\star$]\label{exo:fm:the-risk-management-function:3}
A \omterm{def:fm:the-risk-management-function:cro}{risk committee} meets every fifth business day. How long, on average, does a desk wait for it after a breach opens? And with meetings every
tenth day?
\end{exercise}

\begin{exercise}[$\star\star$]\label{exo:fm:the-risk-management-function:4}
Why is a backdated limit increase a red flag even when the increase itself is justified?
\end{exercise}

\begin{exercise}[$\star\star$]\label{exo:fm:the-risk-management-function:5}
In \cref{tab:fm:the-risk-management-function:regimes}, why do breaches last longer when the committee approves increases, and why does requiring the
desk to cut while it waits bring the duration back down?
\end{exercise}

\begin{exercise}[$\star\star$]\label{exo:fm:the-risk-management-function:6}
Explain why a head of desk approving increases for their own desk satisfies the \omterm{def:fm:the-risk-management-function:foureyes}{four-eyes principle} on paper but not in fact.
\end{exercise}

\begin{exercise}[$\star\star\star$]\label{exo:fm:the-risk-management-function:7}
\emph{Coding.} Run \texttt{fm\_risk.run} for the desk regime and list the breaches of desk 19 before its model change. Which are flagged, and why?
\end{exercise}

\begin{exercise}[$\star\star\star$]\label{exo:fm:the-risk-management-function:8}
\emph{Find the flaw.} ``Our risk function is independent: the \omterm{def:fm:the-risk-management-function:cro}{chief risk officer} reports to the head of trading, who understands the risks better
than anyone.''
\end{exercise}

\section{Problem: Raised to Fit}

\begin{problem}[{Weekend problem --- raised to fit}]\label{pb:fm:the-risk-management-function:1}
A new \omterm{def:fm:the-risk-management-function:cro}{chief risk officer} inherits a firm where heads of desk approve their own temporary increases, and must propose a new \omterm{def:fm:the-risk-management-function:escalation}{escalation procedure}
to the \omterm{def:fm:the-risk-management-function:cro}{risk committee}.

\medskip\noindent\textbf{Part I --- The function.}
\begin{enumerate}
\item Define the \omterm{def:fm:the-risk-management-function:cro}{chief risk officer} and the \omterm{def:fm:the-risk-management-function:cro}{risk committee}.
\item What three things make the \omterm{def:fm:the-risk-management-function:cro}{chief risk officer} independent in practice?
\item What did the Senior Supervisors Group find distinguished the firms that avoided significant losses in 2007?
\item Define the \omterm{def:fm:the-risk-management-function:foureyes}{four-eyes principle} and give two decisions it should cover.
\end{enumerate}

\medskip\noindent\textbf{Part II --- The public record.}
\begin{enumerate}[resume]
\item Give the path of the firm-wide limit in the examiner's report, with the dates.
\item What did the backdating of January 2008 do to the record of breaches?
\item What did the stress tests leave out, and what position was left out of the usage calculation?
\item What did the examiner conclude on the business judgement rule, and what does that mean for a risk function?
\end{enumerate}

\medskip\noindent\textbf{Part III --- The register.}
\begin{enumerate}[resume]
\item Define an \omterm{def:fm:the-risk-management-function:escalation}{escalation procedure} and a \omterm{def:fm:the-risk-management-function:escalation}{risk exception}.
\item State the chapter's severity rules and grade three breaches of your choice.
\item State and prove \cref{prop:fm:the-risk-management-function:wait}.
\item Give the shares of breaches closed by increases under the three regimes.
\item Give the mean duration, breach-days a year and time above the original limit under each.
\item What does the register flag, and what share of breaches is flagged under each regime?
\end{enumerate}

\medskip\noindent\textbf{Part IV --- The proposal.}
\begin{enumerate}[resume]
\item Why is the median duration the same in every regime, and what statistic shows the difference?
\item Describe the model-change episode of \cref{fig:fm:the-risk-management-function:window}.
\item What point-in-time records must the register keep?
\item Define \omterm{def:fm:the-risk-management-function:culture}{risk culture} and name three register metrics that trace it.
\item State the \emph{named result}: the share of breaches resolved by raising the limit, the median time to resolution, and how both move when
increases need the \omterm{def:fm:the-risk-management-function:cro}{risk committee}'s approval.
\item In two sentences, write the new \omterm{def:fm:the-risk-management-function:escalation}{escalation procedure}.
\end{enumerate}
\end{problem}

\section{Interview questions}

\begin{interviewq}[$\star$ \iqroles{risk}]\label{iq:fm:the-risk-management-function:1}
To whom should a \omterm{def:fm:the-risk-management-function:cro}{chief risk officer} report, and why?
\end{interviewq}

\begin{interviewq}[$\star$ \iqroles{trader, risk}]\label{iq:fm:the-risk-management-function:2}
Your desk breaches its VaR limit by 3\% at the close. What happens next?
\end{interviewq}

\begin{interviewq}[$\star\star$ \iqroles{risk}]\label{iq:fm:the-risk-management-function:3}
A desk asks for its third temporary increase in two months. What do you do?
\end{interviewq}

\begin{interviewq}[$\star\star$ \iqroles{risk, researcher}]\label{iq:fm:the-risk-management-function:4}
A desk proposes a new VaR model that lowers its measured risk by 20\% while it is in breach. How do you handle it?
\end{interviewq}

\begin{interviewq}[$\star\star$ \iqroles{risk}]\label{iq:fm:the-risk-management-function:5}
What would you look for in a firm's breach register to judge its \omterm{def:fm:the-risk-management-function:culture}{risk culture}?
\end{interviewq}

\begin{interviewq}[$\star\star\star$ \iqroles{risk, developer}]\label{iq:fm:the-risk-management-function:6}
Design the data model of a breach register that cannot be rewritten after the fact.
\end{interviewq}
